% 钱德拉塞卡拉·拉曼（综述）
% license CCBYSA3
% type Wiki

本文根据 CC-BY-SA 协议转载翻译自维基百科\href{https://en.wikipedia.org/wiki/C._V._Raman}{相关文章}

钱德拉塞卡拉·文卡塔·“C. V.”·拉曼爵士（Sir Chandrasekhara Venkata "C. V." Raman，/ˈrɑːmən/，\(^\text{[2]}\)音译“拉曼”；泰米尔语：சந்திரசேகர வெங்கட ராமன்，罗马化：Cantiracēkara Veṅkaṭa Rāmaṉ；1888年11月7日—1970年11月21日）是一位印度物理学家\(^\text{[3]}\)，以其在光散射领域的研究而闻名。\(^\text{[4]}\)他使用自己研制的分光仪，与学生K. S. 克里希南共同发现：当光线穿过透明物质时，被偏转的光会改变其波长。这种此前未知的光散射现象，他们称之为“修饰散射”，后来被命名为“拉曼效应”或“拉曼散射”。1930年，拉曼因这一发现获得诺贝尔物理学奖，成为第一位获得该奖项的亚洲人和非白人。\(^\text{[5]}\)

拉曼出生于泰米尔婆罗门家庭，自幼聪颖过人。他分别在11岁和13岁时，就完成了圣阿洛伊修斯英印高级中学的中学和高中课程。16岁时，他以优异成绩从马德拉斯大学附属的总督学院毕业，并以物理学荣誉学士的身份名列榜首。他的第一篇研究论文《光的衍射》发表于1906年，当时他仍是一名研究生。次年，他获得了硕士学位。19岁时，拉曼进入加尔各答的印度财政服务部门，担任副会计长。在那里，他结识了印度科学促进会（Indian Association for the Cultivation of Science，简称 IACS）——印度首个科研机构。该机构允许他独立开展研究，也正是在这里，拉曼在声学与光学领域作出了重要贡献。

1917年，拉曼受阿修托什·穆克吉任命，出任加尔各答大学下属拉贾巴扎尔理学院的首任帕利特物理学教授。在他首次前往欧洲的旅途中，亲眼目睹地中海的蓝色使他产生了强烈的好奇心——他认为当时流行的解释（即海洋的蓝色是由于天空中瑞利散射光的反射所致）是错误的。1926年，他创办了《印度物理学杂志》。1933年，他迁往班加罗尔，成为印度科学研究所的首位印度籍所长。同年，他还创立了印度科学院。1948年，他建立了拉曼研究所，并在此工作直至生命的最后时光。

拉曼效应于1928年2月28日被发现。印度政府将这一天定为“全国科学日”，每年举行纪念活动。
\subsection{早年生活与教育}
拉曼出生于英属印度马德拉斯省的蒂鲁吉拉帕利（Tiruchirappalli，今印度泰米尔纳德邦的蒂鲁吉拉帕利），父母均为泰米尔·艾耶尔婆罗门——父亲钱德拉塞卡拉·拉曼纳森·艾耶，\(^\text{[6][7]}\)母亲帕尔瓦蒂·阿玛尔。\(^\text{[8]}\)他是家中八个孩子中的第二个。\(^\text{[9]}\)他的父亲是一名中学教师，收入微薄。拉曼曾幽默地回忆说：“我出生时嘴里含着一把铜勺。因为在我出生时，父亲那‘辉煌’的薪水是每月十卢比！”\(^\text{[10]}\)1892年，全家迁往安得拉邦的维萨卡帕特南（当时称为Vizagapatam或Vizag），因为他的父亲被任命为A.V. 纳拉辛哈·饶夫人学院物理系教师。\(^\text{[11]}\)

拉曼在维萨卡帕特南的圣阿洛伊修斯英印高级中学接受教育。他在11岁时通过了中学毕业考试（matriculation），13岁时以奖学金身份通过了艺术初级考试（First Examination in Arts，相当于今日的大学预科课程），并在两项考试中均名列安得拉邦学校委员会（今安得拉邦中等教育委员会）榜首。
